\section{Introduction}
Industrial anomaly detection (AD) is often deployed in settings where defective samples are rare, diverse, and expensive to annotate. Few-shot AD addresses this constraint by allowing only a small number of normal support images per category, commonly $k=1$--$8$. Recent progress has been driven by frozen vision foundation models: patch features from DINOv2~\cite{dinov22023} can be compared against normal memory banks, as in AnomalyDINO~\cite{anomalydino2024}, or projected onto a normal PCA/subspace model, as in SubspaceAD~\cite{subspacead2026}. These approaches provide strong anomaly ranking without fully supervised defect training.

However, ranking and reliability are different problems. A raw nearest-neighbor distance or subspace residual can produce good AUROC while remaining poorly calibrated: an image with score $0.8$ does not necessarily have an $80\%$ probability of being anomalous. This distinction matters in industrial inspection, where operators need not only a ranked alarm list but also a reliable sense of uncertainty, false-alarm risk, and degradation under shift. Prior work has already shown that DINOv2-based few-shot AD can be poorly calibrated and adversarially fragile~\cite{khan2025calibration}, so improving reliability cannot be treated as a minor post-processing detail.

We therefore study a decoupled design. The ranking score remains the frozen DINOv2 PCA/subspace residual, preserving a low-storage detector and avoiding claims that calibration improves AUROC by changing the anomaly score. Reliability is handled separately by probability calibration and conformal views. Our strongest variant uses leave-one-image-out (LOIO) conformal calibration over the few normal support images: each support image is scored as held out, producing normal nonconformity values from which test p-values are derived. The resulting $1-p$ score is used as a probability-like reliability view, not as a replacement for the raw ranking.

This framing also clarifies our novelty boundary. We do not claim that DINOv2 features, PCA residuals, Platt scaling, conformal prediction, or gated experts are individually new. Instead, our contribution is a protocol-locked empirical and methodological integration: a low-storage few-shot DINOv2 subspace detector with decoupled reliability routing, evaluated across clean, corruption, transfer, and representative routing settings.

Our contributions are:
\begin{itemize}[leftmargin=*]
    \item We propose a low-storage decoupled DINOv2 subspace AD pipeline that keeps PCA residuals for anomaly ranking and uses separate probability/reliability views for decision support.
    \item We introduce LOIO conformal reliability for few-shot normal support and compare it against Vector Platt, Shift-Aware Platt, weighted Platt, and weighted conformal variants.
    \item We provide full VisA corruption evidence showing that LOIO conformal reliability reduces ECE relative to Platt-style calibrators across all tested $k$ and corruption cells while preserving the underlying ranking.
    \item We report clean accuracy-storage, transfer calibration, pixel-level, robustness, and routing diagnostics, including negative results: no MVTec AUROC SOTA claim and no adversarial robustness claim.
\end{itemize}
